\chapter{矩阵特征值}
\label{chap:eigenvalue}

\begin{definition}
    {矩阵特征值问题}{}
    给定$n\times n$矩阵$A$，寻找$\lambda\in\CC$和非零向量$x\in\CC^n$，使得
    \begin{equation}
        Ax=\lambda x,
    \end{equation}
    $\lambda$称为$A$的特征值(eigenvalue)，$x$称为对应的特征向量(eigenvector)。
\end{definition}

\begin{corollary}
    矩阵特征值满足特征多项式方程
    \begin{equation}
        \det(\lambda I-A)=0.
    \end{equation}
\end{corollary}

\paragraph{特征值的分布估计}

由式\eqref{eqn:rho<=norm} $\sigma(A)\leq\norm A$可知，在复平面上，$A$的特征值都位于以原点为中心、(某一)矩阵范数$\norm A$为半径的圆盘内。
当然，这种对特征值界限的估计过于粗糙。

\begin{theorem}
    {Gershgorin圆盘定理·一}{Gershgorin circle theorem I}
    定义$A$的第$i$行Gershgorin圆盘为
    \begin{equation}
        D_i:=\set{z\in\CC}{\abs{z-a_{ii}}\leq r_i},\quad r_i:=\sum_{j\neq i}\abs{a_{ij}}.
    \end{equation}
    则$A$的所有特征值都位于这些圆盘的并集中，即
    \begin{equation}
        \lambda\in\bigcup_{i=1}^nD_i.
    \end{equation}
\end{theorem}

\begin{proof}
    设$\lambda$为$A$的特征值，$x$为对应的特征向量。
    令$k=\argmax_i\abs{x_i}$，则
    \[
        \abs{\lambda - a_{kk}}\abs{x_k}=\biggabs{\sum_{j\neq k}a_{kj}x_j}\leq\sum_{j\neq k}\abs{a_{kj}}\abs{x_j}\leq r_k\abs{x_k}.
    \]
    故$\abs{\lambda - a_{kk}}\leq r_k$，即$\lambda\in D_k$。
\end{proof}

\begin{theorem}
    {Gershgorin圆盘定理·二}{Gershgorin circle theorem II}
    如果$k$个Gershgorin圆盘$D_{i_1},\ldots,D_{i_k}$与其他圆盘互不相交，则它们的并集中恰有$k$个特征值。
\end{theorem}

\begin{proof}
    将$A=L+D+U$分解为严格下三角矩阵$L$、对角矩阵$D$和严格上三角矩阵$U$之和。定义
    \[
        A(\theta):=D+\theta(L+U),\quad\theta\in[0,1].
    \]
    显然$A(1)=A,\,A(0)=D$。
    易得$A(\theta)$的Gershgorin圆盘为
    \[
        D_i(\theta)=\set{z\in\CC}{\abs{z-a_{ii}}\leq \theta r_i},
    \]
    因此$D_i(\theta)\subseteq D_i(1)=D_i$。
    由于$D_{i_1},\ldots,D_{i_k}$与其他圆盘互不相交，故$D_{i_1}(\theta),\ldots,D_{i_k}(\theta)$也与其他圆盘互不相交。

    易得$A(\theta)$的特征多项式关于$\theta$连续变化，故其特征值也关于$\theta$连续变化。
    注意到$D_i(0)$退化为单点$\{a_{ii}\}$，每个$D_i(0)$中都恰有一个特征值。
    考虑$\theta$从0变化到1的过程，由于$D_{i_1}(\theta),\ldots,D_{i_k}(\theta)$与其他圆盘互不相交，其并集中始终有$k$个特征值。
\end{proof}

\begin{center}
    \tikzfigure{gershgorin-disks}
    \captionof{figure}{Gershgorin圆盘定理}
    \label{fig:Gershgorin circle theorem}
\end{center}

\begin{corollary}
    为了更好的估计，可以设法缩小一些Gershgorin圆盘的半径。
    例如$A$与其相似变换$BAB\inv$有相同的特征值，若$B=\diag(b_1,\ldots,b_n)$，则Gershgorin圆盘的半径变为
    \[
        r_i'=\sum_{j\neq i}\biggabs{\frac{b_j}{b_i}a_{ij}}.
    \]
\end{corollary}

\begin{example}
    {}{}
    考虑矩阵$A$及其相似变换$BAB\inv$
    \[
        A=\begin{bmatrix}
            0.9&0.01&0.12\\
            0.01&0.8&0.13\\
            0.01&0.02&0.4
        \end{bmatrix},\quad
        BAB\inv=\begin{bmatrix}
            0.9&0.01&0.012\\
            0.01&0.8&0.013\\
            0.1&0.2&0.4
        \end{bmatrix},
    \]
    其中选取了$B=\diag(1,1,0.1)$，原先$A$的$D_1,D_2$有较大重叠而$D_3$半径较小且孤立，变换后$D_1',D_2'$半径大幅缩小，从而更好地分离了三个圆盘。
\end{example}


